\cleardoublepage % memastikan bab baru dimulai di halaman ganjil (kanan)
\chapter{PENDAHULUAN}
\label{chap:pendahuluan}
% =========================================================
\section{Latar Belakang}
% =========================================================
 
Perkembangan \emph{Large Language Model} (LLM) dalam beberapa tahun terakhir telah mendorong transformasi signifikan pada berbagai aplikasi kecerdasan buatan, termasuk natural language understanding, text generation, conversational agents, machine translation, serta code generation \autocite{kumar2024,qin2025}. Model-model seperti GPT-4 dan LLaMA-2 telah menunjukkan kemampuan luar biasa dalam memahami dan menghasilkan teks natural yang menyerupai tulisan manusia. \textcite{Haque2025} menegaskan bahwa LLM bukan hanya mengubah cara perangkat lunak dikembangkan, tetapi secara fundamental mendefinisikan ulang peran pengembang dalam siklus hidup rekayasa perangkat lunak (\emph{software engineering}). Adopsi LLM di industri perangkat lunak juga meningkat secara signifikan, sebagaimana dikonfirmasi oleh \textcite{Ronanki2026} yang melakukan studi multi-kasus pada tiga organisasi yang menggunakan LLM dalam aktivitas rekayasa perangkat lunak dan mengidentifikasi perlunya praktik terbaik yang dapat ditindaklanjuti untuk adopsi LLM yang efisien dan bertanggung jawab. Meskipun demikian, di balik kemampuan generatif tersebut, terdapat permasalahan mendasar yang belum sepenuhnya terpecahkan yaitu kecenderungan LLM untuk menghasilkan informasi yang tampak meyakinkan namun faktanya tidak akurat atau bahkan sepenuhnya merupakan fabrikasi \autocite{Ji2023}.
 
Fenomena ini dikenal sebagai \emph{hallucination} (halusinasi) dan menjadi salah satu hambatan paling kritis dalam penerapan LLM pada domain yang membutuhkan keakuratan tinggi. \textcite{Ji2023} mendefinisikan halusinasi sebagai kondisi di mana model menghasilkan keluaran yang tidak sesuai dengan fakta yang dapat diverifikasi, tidak konsisten dengan konteks yang diberikan, atau tidak memiliki dasar empiris yang dapat ditelusuri. \textcite{Huang2025} dalam survei komprehensifnya di \emph{ACM Transactions on Information Systems} mengidentifikasi tiga lapisan akar penyebab halusinasi yang saling terkait: (1) kualitas dan kelengkapan data pelatihan yang tidak merata, (2) keterbatasan
proses penyelarasan (\emph{alignment}) model terhadap pengetahuan faktual, dan (3) mekanisme inferensi yang cenderung menghasilkan teks yang terdengar koheren tanpa jaminan akurasi faktual. \textcite{Chang2026} melengkapi analisis ini
dengan mengidentifikasi tiga tantangan fundamental dalam generasi teks berbasis LLM yang memperkuat kecenderungan halusinasi: (1) pemahaman instruksi yang tidak memadai pada konteks panjang dan percakapan multi-gilir, (2) penalaran intensi yang tidak konsisten, khususnya ketika instruksi mengandung informasi ambigu, dan (3) generasi dialog yang tidak andal (\emph{unreliable dialogue generation}) yang menghasilkan konten tidak stabil dari prompt yang serupa.
 
Dalam konteks pemrograman dan rekayasa perangkat lunak, halusinasi LLM bermanifestasi dalam berbagai bentuk yang berpotensi menimbulkan dampak langsung terhadap kualitas sistem. \textcite{Alomari2026} dalam tinjauan literatur sistematis terhadap 49 studi penggunaan LLM untuk peningkatan kualitas kode menemukan bahwa kode yang dihasilkan atau direfaktorisasi oleh LLM tidak dapat diandalkan dan masih menghadapi tantangan besar dalam hal optimasi dan keandalan. Sebuah temuan yang mengonfirmasi bahwa pengetahuan parametrik LLM saja tidak cukup untuk menjamin keakuratan di domain teknis yang terus berkembang. \textcite{DaSilva2025} memperkuat bukti ini melalui studi empiris yang membandingkan jawaban ChatGPT dan LLaMA terhadap pertanyaan Stack Overflow. Hasil penelitian \textcite{DaSilva2025} menunjukkan bahwa meskipun LLM menantang keahlian manusia, LLM tidak melampaui keahlian komunitas manusia untuk sejumlah domain tertentu. Yang lebih signifikan, \textcite{DaSilva2025} mendokumentasikan penurunan nyata aktivitas \emph{posting} pengguna di Stack Overflow sejak hadirnya ChatGPT. Pengguna mulai mengandalkan LLM meskipun keandalannya belum terjamin penuh, sehingga berpotensi memperparah penyebaran informasi teknis yang tidak akurat.
 
Stack Overflow yang merupakan platform utama pengembang untuk pertanyaan seputar pemrograman dan pengembangan perangkat lunak, kini berada di persimpangan antara ancaman dan peluang akibat kehadiran LLM. Di satu sisi, aktivitas penggunanya menurun \autocite{DaSilva2025}. Di sisi lain, komunitas Stack Overflow menyimpan lebih dari 58 juta pasang pertanyaan dan jawaban yang telah melewati mekanisme kurasi ketat berupa sistem \emph{upvote/downvote}, penanda jawaban yang diterima (\emph{accepted answer}), dan sistem reputasi kontributor \autocite{Kabir2024}. Sinyal kepercayaan komunitas ini menjadikan Stack Overflow bukan sekadar repositori teks, melainkan basis pengetahuan tervalidasi secara sosial yang potensial digunakan untuk memberikan landasan pengetahuan bagi jawaban LLM berdasarkan fakta yang telah diverifikasi oleh praktisi domain. \textcite{Nashid2023} menunjukkan bahwa pemanfaatan konten Stack Overflow sebagai sumber pengambilan konteks secara signifikan meningkatkan kualitas jawaban LLM pada tugas-tugas terkait kode, menggarisbawahi potensi besar repositori ini.
 
Salah satu pendekatan yang paling menjanjikan untuk memitigasi halusinasi adalah \emph{Retrieval-Augmented Generation} (RAG) \autocite{Lewis2020}. \textcite{YizhengHuang2026} dalam survei mereka di \emph{ACM Computing Surveys} menjelaskan bahwa RAG menggabungkan teknik pengambilan informasi (\emph{information retrieval}) dengan kemampuan generatif LLM untuk mengatasi keterbatasan pengetahuan statis model, sehingga memungkinkan integrasi informasi eksternal yang dinamis dan terkini pada saat inferensi. \textcite{DelaCruzCabello2025} melalui tinjauan literatur sistematis pada domain deteksi anomali \emph{log} sistem komputer menemukan bahwa RAG secara substansial meningkatkan akurasi dan interpretabilitas jawaban LLM tanpa memerlukan \emph{fine-tuning} yang ekstensif, dan bahwa integrasi RAG dengan LLM menunjukkan kemampuan adaptasi yang kuat terhadap lingkungan operasional yang berubah. Temuan ini mengkonfirmasi potensi RAG sebagai strategi mitigasi halusinasi yang praktis dan skalabel di berbagai domain teknis.
 
Meskipun demikian, implementasi RAG konvensional yang menggunakan pencarian
berbasis kemiripan vektor (\emph{dense retrieval}) terhadap potongan-potongan
teks yang tidak terstruktur memiliki keterbatasan mendasar yaitu pendekatan ini
gagal memanfaatkan relasi semantik kompleks dan hubungan kontekstual antar
konsep yang justru sangat penting dalam domain pemrograman \autocite{YizhengHuang2026}. Sebagai respons atas keterbatasan ini, pendekatan
\emph{Graph-based Retrieval-Augmented Generation} (GraphRAG) diusulkan untuk
mengintegrasikan struktur \emph{knowledge graph} ke dalam proses pengambilan
konteks. \textcite{Zhu2026} dalam survei komprehensif mereka di \emph{ACM
Computing Surveys} mengidentifikasi tiga inovasi kunci GraphRAG dibandingkan
RAG konvensional: (1) representasi pengetahuan berbasis graf yang secara
eksplisit mengkodekan relasi antar entitas dan hierarki domain, (2) teknik
pengambilan berbasis graf yang memungkinkan pengambilan konteks dengan
kemampuan penalaran multi-hop, dan (3) algoritma integrasi pengetahuan yang
sadar-struktur untuk menghasilkan teks yang lebih akurat dan koheren.
\textcite{Pan2024} dalam menegaskan bahwa penggabungan LLM dengan \emph{knowledge graph} membentuk paradigma baru yang secara signifikan meningkatkan kemampuan penalaran berbasis pengetahuan dan mengurangi fabrikasi faktual\emph{(factual hallucination)}.
 
Dari sisi evaluasi, \textcite{Gu2026} dalam survei komprehensif mereka
tentang \emph{LLM-as-a-Judge} menegaskan bahwa mengukur keandalan LLM merupakan tantangan signifikan yang membutuhkan desain dan standardisasi yang cermat. Mereka mengusulkan strategi untuk meningkatkan konsistensi evaluasi dan memitigasi bias, sebuah kerangka yang relevan untuk mengukur efektivitas reduksi halusinasi pada sistem berbasis GraphRAG dan menjadi prasyarat agar kontribusi sebuah \emph{framework} mitigasi halusinasi dapat diklaim secara akademis.
 
Meskipun kemajuan di masing-masing bidang LLM, RAG, GraphRAG, dan penambangan pengetahuan komunitas cukup pesat, terdapat kesenjangan metodologis yang signifikan ketika keempat komponen ini dipertimbangkan secara terpadu. Pertama, penelitian yang secara langsung membandingkan
LLM dengan Stack Overflow seperti \textcite{DaSilva2025} dan \textcite{Kabir2024}, bersifat evaluatif. Mereka mengukur dan mendokumentasikan kesenjangan kualitas antara LLM dan keahlian komunitas, tetapi tidak mengusulkan mekanisme arsitektural untuk menjembatani kesenjangan tersebut. \textit{Kedua}, penelitian yang menggunakan konten Stack Overflow sebagai sumber konteks untuk LLM, seperti \textcite{Nashid2023} hanya menerapkan RAG konvensional berbasis kemiripan vektor dan tidak memanfaatkan struktur relasional pengetahuan komunitas maupun sinyal kepercayaan komunitas (\emph{upvote}, \emph{accepted answer}) sebagai bobot kepercayaan dalam proses pengambilan konteks. \textit{Ketiga}, meskipun pendekatan GraphRAG terbukti superior dibandingkan RAG konvensional pada berbagai domain \autocite{Zhu2026}, serta integrasi \emph{knowledge graph} dengan LLM diyakini sebagai paradigma yang menjanjikan \autocite{Pan2024}, belum ada penelitian yang secara khusus membangun \emph{knowledge graph} dari konten Stack Overflow yang mengkodekan sinyal kepercayaan komunitas untuk digunakan sebagai fondasi GraphRAG pada domain pemrograman. Sementara itu, tantangan halusinasi akibat keterbatasan pemahaman instruksi dan generasi dialog yang tidak andal \autocite{Chang2026, Huang2025} tetap menjadi masalah terbuka yang membutuhkan pendekatan grounding berbasis pengetahuan terverifikasi.
 
Kesenjangan-kesenjangan ini secara kolektif membuka peluang kontribusi yang signifikan: pengembangan sebuah \emph{GraphRAG Framework} yang secara
sistematis (1) membangun \emph{knowledge graph} dari konten Stack Overflow yang mengkodekan entitas teknis, relasi semantik, \emph{dan} sinyal
kepercayaan komunitas sebagai bobot kepercayaan eksplisit; (2) merancang mekanisme \emph{graph-guided retrieval} yang memanfaatkan struktur relasional tersebut untuk mengidentifikasi konteks yang lebih relevan dibandingkan RAG konvensional; (3) mengintegrasikan konteks hasil traversal ke dalam proses inferensi LLM; dan (4) mengukur efektivitas reduksi halusinasi secara kuantitatif menggunakan protokol evaluasi yang terstandarisasi.
 
 
% =========================================================
\section{Rumusan Masalah}
% =========================================================
 
Berdasarkan latar belakang yang telah diuraikan, kondisi ideal yang diharapkan adalah tersedianya mekanisme \emph{grounding} berbasis pengetahuan komunitas yang tervalidasi secara sosial sehingga LLM dapat menghasilkan jawaban teknis yang akurat dan dapat dipertanggungjawabkan. Namun, kondisi aktual yang dijumpai menunjukkan kesenjangan metodologis yang signifikan: penelitian yang menggunakan konten Stack Overflow sebagai konteks LLM masih menggunakan RAG konvensional dan tidak memanfaatkan struktur relasional pengetahuan komunitas \autocite{Nashid2023}; pendekatan GraphRAG yang ada belum mengeksploitasi sumber pengetahuan bervalidasi sosial maupun domain pemrograman secara spesifik \autocite{Zhu2026, Pan2024}; sementara itu, tantangan halusinasi dalam generasi teks LLM pada instruksi yang ambigu dan konteks panjang masih belum terpecahkan \autocite{Chang2026, Huang2025}. Kesenjangan antara kondisi ideal dan kondisi aktual tersebut menjadi landasan permasalahan utama dalam penelitian ini, yang dirumuskan sebagai berikut.
 
\begin{quote}
\textbf{Bagaimana merancang dan mengembangkan GraphRAG \emph{Framework} berbasis pengetahuan komunitas Stack Overflow yang secara terukur mampu
mereduksi halusinasi pada \emph{Large Language Model} di domain pemrograman?}
\end{quote}
 
Untuk menjawab pertanyaan utama tersebut, penelitian ini mengajukan pertanyaan turunan sebagai berikut.
 
\begin{enumerate}
    \item Bagaimana membangun \emph{knowledge graph} yang merepresentasikan
    pengetahuan teknis komunitas Stack Overflow secara terstruktur, termasuk
    entitas (konsep, fungsi, pustaka, \emph{error}), relasi antar entitas,
    serta sinyal kepercayaan komunitas (\emph{upvote}, \emph{accepted answer},
    skor reputasi) sebagai bobot kepercayaan eksplisit pada simpul dan sisi
    graf?
 
    \item Bagaimana merancang mekanisme \emph{graph-guided retrieval} yang
    mampu mengidentifikasi subgraf relevan berdasarkan kueri pengguna,
    sedemikian sehingga konteks yang diekstraksi lebih kaya secara semantik
    dan relasional dibandingkan pengambilan berbasis kemiripan vektor pada
    RAG konvensional?
 
    \item Bagaimana mengintegrasikan konteks hasil traversal \emph{knowledge
    graph} Stack Overflow ke dalam proses inferensi LLM sehingga jawaban yang
    dihasilkan memiliki tingkat akurasi faktual yang lebih tinggi dibandingkan
    LLM tanpa mekanisme \emph{grounding}?
 
    \item Bagaimana merancang protokol evaluasi yang mampu mengukur dan
    memvalidasi efektivitas \emph{framework} yang diusulkan dalam mereduksi
    halusinasi LLM menggunakan metrik kuantitatif yang terstandarisasi,
    serta membandingkan kinerja \emph{framework} terhadap dua kondisi
    \emph{baseline}: LLM murni dan RAG konvensional?
\end{enumerate}
 
 
% =========================================================
\section{Tujuan}
% =========================================================
 
Tujuan penelitian ini dibagi menjadi tujuan umum dan tujuan khusus. Tujuan umum penelitian ini adalah mengembangkan \emph{GraphRAG Framework}
berbasis pengetahuan komunitas Stack Overflow sebagai artefak penelitian desain yang mampu mereduksi halusinasi LLM pada domain pemrograman secara terukur, sekaligus menghasilkan kontribusi pengetahuan berupa prinsip-prinsip desain yang dapat digunakan kembali oleh penelitian serupa di domain teknis lainnya.
 
Adapun tujuan khusus dari penelitian ini adalah sebagai berikut.
 
\begin{enumerate}
    \item Membangun \emph{knowledge graph} dari data Stack Overflow yang
    merepresentasikan entitas teknis, relasi semantik antar entitas, dan
    sinyal kepercayaan komunitas sebagai bobot pada simpul dan sisi graf.
 
    \item Merancang dan mengimplementasikan modul \emph{graph-guided
    retrieval} yang memanfaatkan traversal \emph{knowledge graph} untuk
    mengekstraksi konteks yang relevan secara semantik dan relasional
    berdasarkan kueri pengguna, melampaui keterbatasan RAG konvensional
    berbasis kemiripan vektor.
 
    \item Mengembangkan modul \emph{grounded generation} yang mengintegrasikan
    konteks hasil traversal graf ke dalam mekanisme \emph{prompting} LLM
    untuk menghasilkan jawaban teknis yang memiliki \emph{grounding} faktual
    lebih kuat dibandingkan LLM tanpa mekanisme pengambilan konteks eksternal.
 
    \item Melakukan evaluasi untuk mengukur tingkat
    reduksi halusinasi \emph{framework} yang diusulkan menggunakan metrik
    kuantitatif, serta membandingkan hasilnya terhadap
    dua \emph{baseline} yaitu LLM murni dan RAG konvensional.
 
\end{enumerate}
 
 
% =========================================================
\section{Manfaat Penelitian}
% =========================================================
 
Penelitian ini diharapkan memberikan manfaat pada dua dimensi yang saling melengkapi, yaitu manfaat teoritis dan manfaat praktis.
 
\subsection*{Manfaat Teoritis}
 
\begin{enumerate}

    \item \textbf{Model arsitektur GraphRAG untuk domain \emph{question answering}.}
    \emph{Framework} yang diusulkan memberikan kontribusi berupa model arsitektur
    GraphRAG yang secara eksplisit mempertimbangkan validasi sosial komunitas
    sebagai dimensi kepercayaan dalam proses pengambilan konteks. Model ini
    memperkaya literatur akademik pada persimpangan antara \emph{information
    retrieval}, \emph{knowledge graph}, dan \emph{large language model}.

    \item \textbf{Kerangka evaluasi halusinasi yang terstandarisasi.}
    Protokol evaluasi yang dirancang dalam penelitian ini mencakup metrik akurasi faktual, konsistensi, dan tingkat halusinasi dalam tiga kondisi perbandingan, dapat menjadi referensi metodologis bagi penelitian mitigasi halusinasi LLM di berbagai domain aplikasi.
\end{enumerate}
 
\subsection*{Manfaat Praktis}
 
\begin{enumerate}
    \item \textbf{Peningkatan keandalan bantuan pemrograman berbasis LLM.}
    \emph{GraphRAG Framework} yang dihasilkan dapat diintegrasikan ke dalam
    alat bantu pengembang (\emph{developer tools}) berbasis LLM, seperti
    \emph{AI coding assistant}, untuk menghasilkan saran kode dan penjelasan
    teknis yang lebih akurat dan dapat ditelusuri sumbernya.

    \item \textbf{Panduan adopsi LLM yang bertanggung jawab di industri perangkat lunak.}
    Temuan penelitian ini memberikan bukti empiris dan panduan arsitektural bagi
    organisasi perangkat lunak yang ingin mengadopsi LLM secara efisien dan
    bertanggung jawab, khususnya dalam mengelola risiko penyebaran informasi
    teknis yang tidak akurat.

    \item \textbf{Pemanfaatan kembali aset pengetahuan komunitas Stack Overflow.}
    Penelitian ini menunjukkan cara memanfaatkan lebih dari 58 juta pasang
    pertanyaan dan jawaban Stack Overflow yang telah tervalidasi secara sosial
    sebagai fondasi pengetahuan yang terstruktur, sehingga aset komunitas ini
    tidak hanya menjadi repositori pasif tetapi juga menjadi infrastruktur
    \emph{grounding} aktif untuk sistem AI generatif.
\end{enumerate}
 
 
% =========================================================
\section{Ruang Lingkup dan Batasan Masalah}
% =========================================================
 
Penelitian ini berfokus pada perancangan dan pengembangan \emph{GraphRAG Framework} untuk domain pemrograman dengan menggunakan data dari komunitas Stack Overflow sebagai sumber pengetahuan utama. Evaluasi dilakukan pada pertanyaan dan jawaban teknis yang mencakup topik-topik pemrograman umum (\emph{general programming}) seperti Python, JavaScript, dan Java, yang merupakan \emph{tag} dengan volume konten tertinggi di Stack Overflow.
 
\textbf{Batasan masalah:}
\begin{enumerate}
    \item \emph{Framework} yang dikembangkan difokuskan pada domain
    pemrograman dan tidak dimaksudkan untuk langsung digeneralisasi ke domain
    non-teknis tanpa adaptasi pada skema \emph{knowledge graph} dan protokol
    evaluasinya.
    \item Pembangunan \emph{knowledge graph} dibatasi pada ekstraksi entitas
    dan relasi dari teks pertanyaan, jawaban, dan komentar Stack Overflow,
    tanpa menyertakan eksekusi atau validasi kode secara dinamis.
    \item Penelitian ini tidak membahas aspek \emph{fine-tuning} LLM.
    Seluruh eksperimen menggunakan LLM dalam mode \emph{inference} tanpa
    modifikasi parameter model.
    \item Pengukuran halusinasi dibatasi pada dimensi akurasi faktual dan
    konsistensi informasi, tidak mencakup dimensi bias, toksisitas, atau
    kesesuaian gaya bahasa.
    \item Penelitian ini tidak mencakup aspek skalabilitas dan
    \emph{deployment framework} ke lingkungan produksi. Seluruh eksperimen
    dilakukan dalam lingkungan penelitian yang terkontrol.
\end{enumerate}
 
 
% =========================================================
\section{Hipotesis}
% =========================================================
 
Hipotesis penelitian ini diturunkan dari dua premis yang diidentifikasi pada latar belakang dan rumusan masalah.
 
\begin{enumerate}
    \item \textbf{Premis 1.} LLM yang beroperasi hanya mengandalkan
    pengetahuan parametrik menghasilkan tingkat halusinasi yang signifikan
    pada pertanyaan teknis domain pemrograman, karena pengetahuan yang
    dikodekan selama pelatihan bersifat statis dan tidak mencerminkan
    perkembangan teknologi, perubahan API, maupun solusi terkini yang
    didokumentasikan oleh komunitas praktisi \autocite{Ji2023, Huang2025}.
    Hal ini dikonfirmasi oleh \textcite{DaSilva2025} secara empiris dan oleh
    \textcite{Chang2026} yang mengidentifikasi defisiensi fundamental LLM
    dalam menghasilkan dialog yang andal. RAG konvensional mengurangi
    sebagian halusinasi ini, namun masih terbatas karena pengambilan berbasis
    kemiripan vektor gagal memanfaatkan relasi semantik kompleks antar konsep
    teknis \autocite{YizhengHuang2026}.
 
    \item \textbf{Premis 2.} \emph{Knowledge graph} yang dibangun dari
    konten Stack Overflow merepresentasikan pengetahuan teknis dalam struktur
    relasional yang kaya, dan sinyal kepercayaan komunitas (skor \emph{upvote},
    status \emph{accepted answer}) merupakan proksi yang valid untuk kualitas
    dan kebenaran informasi \autocite{DaSilva2025, Kabir2024}. Mekanisme
    \emph{graph-guided retrieval} yang memanfaatkan struktur relasional ini
    mampu mengidentifikasi konteks yang lebih relevan dan terverifikasi
    dibandingkan pencarian vektor sederhana, sebagaimana diindikasikan oleh
    superioritas pendekatan GraphRAG yang dikonfirmasi oleh \textcite{Zhu2026}
    dan paradigma integrasi LLM-\emph{knowledge graph} yang diuraikan oleh
    \textcite{Pan2024}.
\end{enumerate}
 
Berdasarkan kedua premis tersebut, hipotesis penelitian ini dirumuskan
sebagai berikut: \emph{GraphRAG Framework} yang mengintegrasikan
\emph{knowledge graph} berbasis pengetahuan komunitas Stack Overflow dengan
mekanisme \emph{graph-guided retrieval} dan sinyal kepercayaan komunitas
sebagai bobot menghasilkan reduksi halusinasi LLM yang secara statistik
signifikan dibandingkan dua kondisi \emph{baseline} yaitu LLM murni tanpa
mekanisme pengambilan konteks dan LLM dengan RAG konvensional berbasis
\emph{dense retrieval} yang diukur menggunakan metrik akurasi faktual dan
skor konsistensi pada \emph{benchmark} pertanyaan teknis domain pemrograman.
 
 
% =========================================================
\section{Metodologi Penelitian}
% =========================================================
 
Penelitian ini menggunakan kerangka \emph{Design Science Research Methodology} (DSRM) sebagai pendekatan sistematis untuk merancang, mengembangkan, dan mengevaluasi artefak berbasis teknologi \autocite{Peffers2007}. DSRM dipilih karena penelitian ini tidak hanya bertujuan melakukan analisis terhadap fenomena yang ada, tetapi juga menghasilkan artefak nyata berupa \emph{GraphRAG Framework} yang dapat digunakan sebagai solusi terhadap permasalahan halusinasi LLM. Penelitian ini dilaksanakan melalui enam tahapan berikut.
 
\begin{enumerate}
 
\item \textbf{Identifikasi Masalah dan Motivasi.}
Tahap ini dilaksanakan melalui studi literatur sistematis dan analisis
\emph{research gap} untuk memetakan keterbatasan pendekatan yang ada serta
mengidentifikasi peluang integrasi pengetahuan komunitas tervalidasi ke dalam
arsitektur GraphRAG. Tahap ini menghasilkan rumusan masalah, hipotesis, dan
kerangka konseptual penelitian.
 
\item \textbf{Pendefinisian Tujuan Solusi.}
Berdasarkan identifikasi masalah, ditetapkan tujuan solusi berupa spesifikasi
fungsional dan non-fungsional \emph{GraphRAG Framework}: skema \emph{knowledge
graph}, mekanisme \emph{retrieval} berbasis graf, modul \emph{grounded
generation}, dan protokol evaluasi halusinasi yang terstandarisasi.
 
\item \textbf{Perancangan dan Pengembangan Artefak.}
Artefak dirancang sebagai \emph{GraphRAG Framework} dengan empat modul
terintegrasi: (1) modul konstruksi \emph{knowledge graph} dari
\emph{Stack Exchange Data Dump} yang mengkodekan entitas teknis, relasi
semantik, dan sinyal kepercayaan komunitas sebagai bobot; (2) modul
\emph{graph-guided retrieval} berbasis traversal berbobot; (3) modul
\emph{grounded generation} yang menyusun \emph{structured prompt} dari
subgraf hasil ekstraksi; serta (4) modul evaluasi. Seluruh komponen
didokumentasikan menggunakan diagram arsitektur dan spesifikasi antarmuka
antar modul.
 
\item \textbf{Implementasi.}
Implementasi dilakukan pada tumpukan teknologi yang dipilih berdasarkan
kesesuaian dengan kebutuhan pembangunan \emph{knowledge graph} dan integrasi
LLM. Data primer yang digunakan adalah \emph{Stack Exchange Data Dump}
yang difilter berdasarkan rentang waktu 2020--2024 dan \emph{tag} pemrograman
yang relevan.
 
\item \textbf{Evaluasi.}
Evaluasi artefak dilakukan pada dua level: evaluasi kualitas
\emph{knowledge graph} dan evaluasi efektivitas reduksi halusinasi pada
tiga kondisi perbandingan (LLM murni, RAG konvensional, dan \emph{GraphRAG
Framework} yang diusulkan). Pengujian hipotesis dilakukan menggunakan uji
statistik yang sesuai untuk menilai signifikansi perbedaan antar kondisi.
 
\item \textbf{Penarikan Kesimpulan dan Komunikasi Hasil.}
Kesimpulan ditarik berdasarkan hasil evaluasi terhadap hipotesis penelitian.
Hasil penelitian dikomunikasikan dalam bentuk laporan tesis, dokumentasi
teknis \emph{framework}, serta artikel ilmiah yang direncanakan
untuk dipublikasikan pada konferensi atau jurnal bereputasi di bidang
\emph{Natural Language Processing}, \emph{Information Retrieval}, atau
\emph{Software Engineering}.
 
\end{enumerate}
 
 
% =========================================================
\section{Sistematika Penulisan}
% =========================================================
 
Bagian ini menjelaskan gambaran umum isi dari setiap bab dalam laporan
tesis. Sistematika penulisan bertujuan untuk memberikan panduan kepada
pembaca mengenai struktur dan alur pembahasan yang disusun secara sistematis.
Adapun sistematika penulisan dalam laporan tesis ini adalah sebagai berikut.
 
\begin{enumerate}
 
\item \textbf{Bab I Pendahuluan}\\
Berisi latar belakang penelitian, rumusan masalah, tujuan penelitian,
manfaat penelitian, ruang lingkup dan batasan masalah, hipotesis, metodologi
penelitian, serta sistematika penulisan laporan tesis.
 
\item \textbf{Bab II Studi Literatur}\\
Berisi kajian pustaka yang relevan dengan topik tesis, meliputi konsep
dasar LLM dan permasalahan halusinasi, pendekatan RAG dan GraphRAG,
integrasi \emph{knowledge graph} dengan LLM, pemanfaatan pengetahuan
komunitas Stack Overflow dalam sistem AI, serta metodologi evaluasi
keandalan LLM. Kajian ini juga menyajikan analisis \emph{research gap}
yang menjadi landasan \emph{novelty} penelitian.
 
\item \textbf{Bab III Analisis dan Desain}\\
Berisi analisis permasalahan halusinasi LLM secara mendalam, identifikasi
kebutuhan komponen \emph{framework}, pemetaan sumber data Stack Overflow,
analisis karakteristik \emph{knowledge graph} yang dibutuhkan, serta
perbandingan alternatif solusi yang menjadi dasar perancangan
\emph{GraphRAG Framework}.
 
\item \textbf{Bab IV Perancangan}\\
Berisi perancangan solusi yang diusulkan secara rinci, meliputi arsitektur
keseluruhan \emph{GraphRAG Framework}, skema \emph{knowledge graph}
(definisi entitas, relasi, dan atribut bobot kepercayaan), desain
mekanisme \emph{graph-guided retrieval}, desain modul \emph{grounded
generation}, serta protokol evaluasi halusinasi.
 
\item \textbf{Bab V Implementasi}\\
Berisi proses implementasi \emph{GraphRAG Framework} yang telah dirancang,
termasuk pembangunan \emph{knowledge graph} dari \emph{Stack Exchange Data
Dump}, implementasi modul \emph{retrieval} dan \emph{generation}, teknologi
dan pustaka yang digunakan, serta hasil implementasi per modul.
 
\item \textbf{Bab VI Evaluasi}\\
Berisi proses pengujian dan evaluasi terhadap \emph{framework} yang
diimplementasikan, meliputi desain eksperimen, konstruksi \emph{benchmark
dataset}, skenario pengujian tiga kondisi (LLM murni, RAG konvensional,
GraphRAG), hasil pengujian kuantitatif, analisis statistik, serta diskusi
temuan terhadap hipotesis penelitian.
 
\item \textbf{Bab VII Kesimpulan dan Saran}\\
Berisi kesimpulan dari hasil penelitian yang telah dilakukan, prinsip-prinsip
desain (\emph{design principles}) yang dirumuskan, keterbatasan penelitian,
serta saran untuk pengembangan penelitian selanjutnya.
 
\item \textbf{Daftar Pustaka}\\
Berisi seluruh referensi yang digunakan dalam penyusunan laporan tesis.
 
\item \textbf{Lampiran}\\
Berisi dokumen pendukung seperti kode program, skema lengkap \emph{knowledge
graph}, data hasil pengujian secara lengkap, serta dokumen lain yang relevan.
 
\end{enumerate}